% अभ्यास-संग्रह — अध्याय 2
% अध्याय संख्या इस फ़ाइल के नाम से निर्धारित होती है।

\begin{exo}[ठंडी बर्फ़ से उबलते पानी तक: परिमाण-क्रम]{chauffage-glace-eau-ebullition}
\seoready{true}
\keywords{ऊष्मामिति, अवस्था परिवर्तन, गलन, वाष्पीकरण, गुप्त ऊष्मा, परिमाण-क्रम}

\textbf{संख्यात्मक आँकड़े (वायुमंडलीय दाब पर):}
\[
\begin{aligned}
c_{\text{ice}} &= 2{,}1\times10^3\,\text{J}\!\cdot\!\text{kg}^{-1}\!\cdot\!\text{K}^{-1},
& L_f &= 3{,}34\times10^5\,\text{J}\!\cdot\!\text{kg}^{-1},\\
c_{\text{water}} &= 4{,}18\times10^3\,\text{J}\!\cdot\!\text{kg}^{-1}\!\cdot\!\text{K}^{-1},
& L_v &= 2{,}26\times10^6\,\text{J}\!\cdot\!\text{kg}^{-1}.
\end{aligned}
\]
वायुमंडलीय दाब पर, प्रारम्भ में $-100\,{}^\circ\text{C}$ तापमान वाली एक किलोग्राम बर्फ़ को क्रमशः गर्म किया जाता है।

\begin{enumerate}
  \item बर्फ़ को $-100\,{}^\circ\text{C}$ से $0\,{}^\circ\text{C}$ तक गर्म करने के लिए आवश्यक ऊर्जा निकालिए।
  \item $0\,{}^\circ\text{C}$ पर पूरी बर्फ़ को गलाने के लिए आवश्यक ऊर्जा निकालिए।
  \item द्रव पानी को $0\,{}^\circ\text{C}$ से $100\,{}^\circ\text{C}$ तक गर्म करने के लिए आवश्यक ऊर्जा निकालिए।
  \item $100\,{}^\circ\text{C}$ पर पूरे पानी को वाष्पित करने के लिए आवश्यक अतिरिक्त ऊर्जा निकालिए।
  \item किस चरण में सबसे अधिक ऊर्जा लगती है? एक किलोग्राम पानी को $0\,{}^\circ\text{C}$ से $100\,{}^\circ\text{C}$ तक गर्म करने की ऊर्जा की तुलना दस मीटर गिरने वाले द्रव्यमान $m$ की गुरुत्वीय स्थितिज ऊर्जा से कीजिए। समान ऊर्जा मुक्त करने के लिए कितना द्रव्यमान गिरना चाहिए?
  \item अब उलटी प्रक्रिया पर विचार कीजिए। जलवाष्प का द्रव्यमान $m_v$ कार के सामने के शीशे पर संघनित होता है और बने हुए द्रव पानी को बाद में ठंडा नहीं माना जाता। संघनन में मुक्त ऊष्मा $Q_{\text{rel}}$ लिखिए और $m_v=1{,}0\times10^{-2}\,\text{kg}$ के लिए निकालिए। शीशे के तापमान पर $L_v=2{,}45\times10^6\,\text{J}\!\cdot\!\text{kg}^{-1}$ लीजिए।
\end{enumerate}

\begin{indication}
अवस्था परिवर्तन के बिना गर्म करने के लिए $Q=mc\Delta T$ और स्थिर तापमान पर अवस्था परिवर्तन के लिए $Q=mL$ प्रयोग कीजिए। ऊँचाई $h$ पर द्रव्यमान $m$ की गुरुत्वीय स्थितिज ऊर्जा $E_p=mgh$ है; $g=9{,}81\,\text{m}\!\cdot\!\text{s}^{-2}$ लीजिए।
\end{indication}

\begin{solution}
\textbf{1.} बर्फ़ को गलन तापमान तक गर्म करने के लिए
\[
Q_1=mc_{\text{ice}}\Delta T
=1{,}0\times2{,}1\times10^3\times100=2{,}1\times10^5\,\text{J}.
\]
\textbf{2.} गलन के दौरान तापमान $0\,{}^\circ\text{C}$ रहता है:
\[
Q_2=mL_f=1{,}0\times3{,}34\times10^5=3{,}34\times10^5\,\text{J}.
\]
\textbf{3.} द्रव पानी को गर्म करने के लिए
\[
Q_3=mc_{\text{water}}\Delta T
=1{,}0\times4{,}18\times10^3\times100=4{,}18\times10^5\,\text{J}.
\]
\textbf{4.} पूर्ण वाष्पीकरण के लिए अतिरिक्त
\[
Q_4=mL_v=1{,}0\times2{,}26\times10^6=2{,}26\times10^6\,\text{J}
\]
चाहिए। केवल यह चरण ही कुछ मेगाजूल ऊर्जा माँगता है।

\textbf{5.} वाष्पीकरण स्पष्ट रूप से सबसे अधिक ऊर्जा-खपत वाला चरण है:
\[
Q_4=2{,}26\times10^6>Q_3=4{,}18\times10^5>Q_2=3{,}34\times10^5>Q_1=2{,}1\times10^5\,\text{J}.
\]
वाष्पीकरण में द्रव पानी को $0$ से $100\,{}^\circ\text{C}$ तक गर्म करने की अपेक्षा लगभग $(2{,}26\times10^6)/(4{,}18\times10^5)\approx5{,}4$ गुना ऊर्जा लगती है।

$h=10\,\text{m}$ से गिरते द्रव्यमान $m$ के लिए $E_p=mgh$। $mgh=Q_3$ रखने पर
\[
m=\frac{Q_3}{gh}=\frac{4{,}18\times10^5}{9{,}81\times10}
\approx4{,}26\times10^3\,\text{kg}.
\]
अतः एक किलोग्राम पानी को $0$ से $100\,{}^\circ\text{C}$ तक गर्म करने जितनी ऊर्जा मुक्त करने के लिए लगभग $4{,}3$ टन को दस मीटर गिराना होगा। यह बहुत बड़ी मात्रा है और बताती है कि घरेलू ऊर्जा उपयोग में पानी गर्म करने का भाग महत्त्वपूर्ण क्यों है। पानी की विशिष्ट ऊष्मा क्षमता सामान्य धातुओं से भी बहुत बड़ी है; उदाहरण के लिए यह ताँबे की लगभग दस गुनी है (आगे के अभ्यास देखें)।

\textbf{6.} संघनन वाष्पीकरण की उलटी प्रक्रिया है; वाष्प शीशे को गुप्त ऊष्मा देता है:
\[
Q_{\text{rel}}=m_vL_v.
\]
$m_v=1{,}0\times10^{-2}\,\text{kg}$ के लिए
\[
Q_{\text{rel}}=1{,}0\times10^{-2}\times2{,}45\times10^6
=2{,}45\times10^4\,\text{J}.
\]
इसलिए थोड़ी-सी वाष्प का संघनन भी उपेक्षणीय न होने वाली ऊष्मा मुक्त कर सकता है, जिसे शीशा और उसका परिवेश ग्रहण करते हैं।
\end{solution}
\end{exo}

\begin{exo}[एक बड़े पेड़ का वाष्पोत्सर्जन: प्राकृतिक वातानुकूलक?]{odg-evapotranspiration-arbre}
\seoready{true}
\keywords{परिमाण-क्रम, वाष्पोत्सर्जन, पेड़, गुप्त ऊष्मा, शीतलन शक्ति, वातानुकूलक, COP}

गर्म और सूखे दिन में पर्याप्त पानी पाने वाला एक बड़ा पेड़ वाष्पोत्सर्जन से लगभग $500\,\text{L}=5{,}0\times10^{-1}\,\text{m}^3$ पानी छोड़ सकता है। वाष्पोत्सर्जन मुख्यतः प्रकाश में होता है, इसलिए मानते हैं कि यह बारह घंटे चलता है। पेड़ की छतरी भूमि पर $A_{\text{tree}}=50\,\text{m}^2$ क्षेत्र ढकती है। पानी का घनत्व और वायुमंडलीय दाब पर वाष्पीकरण की गुप्त ऊष्मा हैं
\[
\rho_{\text{water}}=1{,}0\times10^3\,\text{kg}\!\cdot\!\text{m}^{-3},
\qquad L_v=2{,}45\times10^6\,\text{J}\!\cdot\!\text{kg}^{-1}.
\]
\begin{enumerate}
  \item दिन भर में वाष्पोत्सर्जित पानी का द्रव्यमान और उसे वाष्पित करने की ऊर्जा निकालिए।
  \item समझाइए कि इससे शीतलन क्यों होता है।
  \item सक्रिय वाष्पोत्सर्जन के बारह घंटों में पेड़ की प्रति वर्ग मीटर औसत शीतलन शक्ति निकालिए।
  \item तुलना के लिए $P_{\mathrm{el}}=2{,}0\times10^3\,\text{W}$ विद्युत शक्ति और $\mathrm{COP}=3{,}0$ वाले वातानुकूलक को लीजिए, जो $A_{\text{room}}=20\,\text{m}^2$ कमरे को ठंडा करता है। परिभाषा से $\mathrm{COP}=P_{\text{cool}}/P_{\mathrm{el}}$। शीतलन शक्ति और प्रति वर्ग मीटर शक्ति निकालकर पेड़ से तुलना कीजिए।
  \item प्रति इकाई क्षेत्रफल की शक्तियाँ समान परिमाण-क्रम की होने पर भी पेड़ और वातानुकूलक समान अनुभूत शीतलन देते हैं, ऐसा निष्कर्ष क्यों नहीं निकाला जा सकता?
  \item क्या बहुत-से घरेलू पौधे घर को ठंडा कर सकते हैं? (सामान्य ज्ञान: उत्तर जैविक प्रक्रियाओं पर निर्भर है।)
\end{enumerate}
\begin{indication}
$m=\rho V$, $Q_{\mathrm{evap}}=mL_v$ और $P=Q/\tau$ प्रयोग कीजिए। वातानुकूलक के लिए $P_{\text{cool}}=\mathrm{COP}\,P_{\mathrm{el}}$।
\end{indication}
\begin{solution}
\textbf{1.} पानी का द्रव्यमान
\[
m=\rho_{\text{water}}V=1{,}0\times10^3\times5{,}0\times10^{-1}
=5{,}0\times10^2\,\text{kg}.
\]
वाष्पीकरण ऊर्जा
\[
Q_{\mathrm{evap}}=mL_v=5{,}0\times10^2\times2{,}45\times10^6
=1{,}225\times10^9\,\text{J}.
\]
इसका परिमाण-क्रम एक गीगाजूल है।

\textbf{2.} वाष्पीकरण ऊष्माशोषी अवस्था परिवर्तन है। ऊर्जा पत्तियों, मिट्टी और आसपास की हवा से ली जाती है, जिससे पत्तियाँ और निकट का वातावरण ठंडे होते हैं।

\textbf{3.} बारह घंटे $\tau=12\times3600=4{,}32\times10^4\,\text{s}$ हैं। अतः
\[
P_{\text{tree}}=\frac{Q_{\mathrm{evap}}}{\tau}
=\frac{1{,}225\times10^9}{4{,}32\times10^4}\approx2{,}84\times10^4\,\text{W},
\]
और प्रति इकाई क्षेत्रफल
\[
\frac{P_{\text{tree}}}{A_{\text{tree}}}
=\frac{2{,}84\times10^4}{50}\approx5{,}7\times10^2\,\text{W}\!\cdot\!\text{m}^{-2}.
\]
\textbf{4.} वातानुकूलक की उपयोगी शीतलन शक्ति
\[
P_{\text{cool}}=\mathrm{COP}\,P_{\mathrm{el}}
=3{,}0\times2{,}0\times10^3=6{,}0\times10^3\,\text{W}.
\]
कमरे के प्रति इकाई क्षेत्रफल पर
\[
\frac{P_{\text{cool}}}{A_{\text{room}}}
=\frac{6{,}0\times10^3}{20}=3{,}0\times10^2\,\text{W}\!\cdot\!\text{m}^{-2}.
\]
इस मॉडल में पेड़ की प्रति क्षेत्रफल वाष्पोत्सर्जन शक्ति सामान्य वातानुकूलक की लगभग दुगुनी है।

\textbf{5.} तुलना केवल ऊर्जा प्रवाह की है। वातानुकूलक बंद और नियंत्रित आयतन ठंडा करता है, जबकि पेड़ की वाष्प वायुमंडल में फैलती है और हवा गर्म वायु लाती है। अनुभूति वेंटिलेशन, आर्द्रता, तापमान और पेड़ से दूरी पर निर्भर करती है। वाष्पोत्सर्जन धूप, हवा, आर्द्रता, प्रजाति और उपलब्ध पानी से बदलता है तथा सूखे में बहुत घट सकता है। पेड़ छाया भी देता है, जिससे भूमि और लोगों को मिलने वाला सौर विकिरण सीधे घटता है। इसलिए गणना परिमाण-क्रमों की तुलना है, समान तापीय आराम का प्रमाण नहीं।

\textbf{6.} व्यवहार में घरेलू पौधों का शीतलन बहुत कम होगा। अधिकांश पौधों में प्रकाश रंध्रों को खोलता है; इन छोटे पत्ती-छिद्रों से जलवाष्प निकलती है। घर के भीतर प्रकाश प्रायः कम होता है, इसलिए रंध्र कम खुलते हैं और वाष्पोत्सर्जन घटता है। कम वेंटिलेशन वाले कमरे में बढ़ती आर्द्रता भी वाष्पीकरण को धीमा करती है।

बहुत-से पौधे थोड़ा स्थानीय वाष्पीकरणीय शीतलन दे सकते हैं, पर वातानुकूलक का स्थान नहीं लेते। ऊर्जा को लगातार घर से निकालने के लिए नम हवा बाहर भेजनी होगी। तेज कृत्रिम प्रकाश वाष्पोत्सर्जन बढ़ा सकता है, लेकिन उसकी विद्युत ऊर्जा लगभग पूरी तरह कमरे में ऊष्मा बनेगी। शुष्क परिवेश के अनुकूल CAM पौधे, जिनमें कैक्टस शामिल हैं, अपवाद हैं: उनके रंध्र मुख्यतः रात में खुलते हैं।
\end{solution}
\end{exo}

\begin{exo}[पानी के दो द्रव्यमानों का मिश्रण]{calorimetrie-melange-eau}
\seoready{true}
\keywords{ऊष्मामिति, मिश्रण, ऊष्मा क्षमता, विशिष्ट ऊष्मा, जोज़ेफ़ ब्लैक, संतुलन तापमान}

$T_1=80\,^\circ\text{C}$ पर पानी का $m_1=0{,}200\,\text{kg}$ भाग ऐसे पात्र में डाला जाता है जिसमें $T_2=15\,^\circ\text{C}$ पर $m_2=0{,}300\,\text{kg}$ पानी है। पात्र आदर्श ऊष्मामापी है: बाहर की ओर ऊष्मा हानि और पात्र की अपनी ऊष्मा क्षमता उपेक्षित हैं। $Q=mc\Delta T$, जहाँ द्रव पानी की विशिष्ट ऊष्मा $c=4{,}18\times10^3\,\text{J}\!\cdot\!\text{kg}^{-1}\!\cdot\!\text{K}^{-1}$ है।
\begin{enumerate}
  \item समझाइए कि गर्म पानी की छोड़ी ऊष्मा ठंडे पानी को पूरी मिलती है और $m_1$, $m_2$, $c$, $T_1$, $T_2$, संतुलन तापमान $T_{eq}$ वाला संरक्षण समीकरण लिखिए।
  \item $T_{eq}$ का व्यंजक और संख्यात्मक मान निकालिए।
  \item मिश्रण में गर्म पानी की छोड़ी ऊष्मा $Q$ निकालिए।
  \item क्या एक ही पदार्थ के ऐसे मिश्रण से $c$ निर्धारित हो सकता है? कारण दीजिए।
\end{enumerate}
\begin{indication}
ऊष्मामापी पृथक और दृढ़ है, इसलिए कुल आंतरिक ऊर्जा $U_1+U_2$ संरक्षित है: एक भाग की छोड़ी ऊष्मा दूसरे को विपरीत चिह्न के साथ मिलती है।
\end{indication}
\begin{solution}
\textbf{1.}
\[
m_1c(T_{eq}-T_1)+m_2c(T_{eq}-T_2)=0.
\]
\textbf{2.} समान गुणक $c$ कट जाता है:
\[
T_{eq}=\frac{m_1T_1+m_2T_2}{m_1+m_2}
=\frac{0{,}200\times80+0{,}300\times15}{0{,}200+0{,}300}=41\,^\circ\text{C}.
\]
\textbf{3.}
\[
Q=m_1c(T_1-T_{eq})=0{,}200\times4{,}18\times10^3\times(80-41)
\approx3{,}26\times10^4\,\text{J}.
\]
ठंडा पानी उतनी ही ऊष्मा ग्रहण करता है:
\[
m_2c(T_{eq}-T_2)=0{,}300\times4{,}18\times10^3\times26
\approx3{,}26\times10^4\,\text{J}.
\]
\textbf{4.} नहीं। एक ही पदार्थ के लिए $c$ ऊर्जा संतुलन के दोनों पदों में आता है और कट जाता है। संतुलन तापमान $c$ पर निर्भर नहीं, बल्कि प्रारम्भिक तापमानों का द्रव्यमान-भारित औसत है। मिश्रण विधि से विशिष्ट ऊष्मा मापने के लिए ज्ञात ऊष्मा क्षमता वाला पदार्थ चाहिए, जैसा अगले अभ्यास में है।
\end{solution}
\end{exo}

\begin{exo}[मिश्रण विधि से धातु की विशिष्ट ऊष्मा का निर्धारण]{calorimetrie-methode-melanges}
\seoready{true}
\keywords{ऊष्मामिति, मिश्रण विधि, विशिष्ट ऊष्मा, ऊष्मामापी, जल तुल्यांक}

अठारहवीं शताब्दी में जोज़ेफ़ ब्लैक की तरह \emph{मिश्रण विधि} से अज्ञात धातु की विशिष्ट ऊष्मा $c_m$ मापनी है। $m=0{,}150\,\text{kg}$ नमूने को $T_m=95\,^\circ\text{C}$ तक गर्म करके शीघ्र ही $m_e=0{,}250\,\text{kg}$ पानी वाले ऊष्मामापी में डालते हैं। पानी के लिए $c_e=4{,}18\times10^3\,\text{J}\!\cdot\!\text{kg}^{-1}\!\cdot\!\text{K}^{-1}$ और $T_e=18\,^\circ\text{C}$ हैं; संतुलन तापमान $T_{eq}=22\,^\circ\text{C}$ है। पहले पात्र, विलोड़क और तापमापी की ऊष्मा क्षमता उपेक्षित करें।
\begin{enumerate}
  \item धातु और पानी के पृथक तंत्र का ऊर्जा संरक्षण $c_m$ को एकमात्र अज्ञात रखते हुए लिखिए।
  \item $c_m$ का व्यंजक और मान निकालिए। यह किस सामान्य धातु के लगभग बराबर है? संदर्भ मानों की इकाई $\text{J}\!\cdot\!\text{kg}^{-1}\!\cdot\!\text{K}^{-1}$ है: ऐलुमिनियम $9{,}0\times10^2$, लोहा $4{,}5\times10^2$, ताँबा $3{,}9\times10^2$, सीसा $1{,}3\times10^2$।
  \item पात्र भी कुछ ऊष्मा ग्रहण करता है। इसे \emph{जल तुल्यांक} $\mu=1{,}5\times10^{-2}\,\text{kg}$ से मॉडल करें: ऊष्मामापी अतिरिक्त पानी द्रव्यमान $\mu$ जैसा है। $c_m$ फिर निकालकर सुधार की दिशा समझाइए।
  \item नमूना शीघ्र डालना और ऊष्मामापी को परिवेशी हवा से अच्छी तरह पृथक करना क्यों आवश्यक है?
\end{enumerate}
\begin{indication}
धातु $Q_m=mc_m(T_m-T_{eq})$ छोड़ती है, जिसे पानी और प्रश्न 3 में ऊष्मामापी पूरा ग्रहण करते हैं: $Q_m=(m_e+\mu)c_e(T_{eq}-T_e)$।
\end{indication}
\begin{solution}
\textbf{1.}
\[
mc_m(T_m-T_{eq})=m_ec_e(T_{eq}-T_e).
\]
\textbf{2.}
\[
c_m=\frac{m_ec_e(T_{eq}-T_e)}{m(T_m-T_{eq})}
=\frac{0{,}250\times4{,}18\times10^3\times(22-18)}{0{,}150\times(95-22)}
\approx3{,}82\times10^2\,\text{J}\!\cdot\!\text{kg}^{-1}\!\cdot\!\text{K}^{-1}.
\]
यह मान ताँबे के बहुत निकट है।

\textbf{3.}
\[
c_m=\frac{(m_e+\mu)c_e(T_{eq}-T_e)}{m(T_m-T_{eq})}
=\frac{(0{,}250+0{,}015)\times4{,}18\times10^3\times4}{0{,}150\times73}
\approx4{,}05\times10^2\,\text{J}\!\cdot\!\text{kg}^{-1}\!\cdot\!\text{K}^{-1}.
\]
सुधार $c_m$ बढ़ाता है: पात्र को उपेक्षित करने से धातु की वास्तव में छोड़ी ऊष्मा और उसकी विशिष्ट ऊष्मा कम आँकी जाती थी।

\textbf{4.} विधि पूरे मापन के दौरान पृथक तंत्र मानती है। नमूना शीघ्र डालने से पानी तक पहुँचने से पहले हवा से ऊष्मा विनिमय घटता है; अच्छा पृथक्करण संतुलन बनने के दौरान हानि घटाता है। कोई भी अनगिना ऊष्मा प्रवाह संतुलन और निकले $c_m$ को बिगाड़ेगा। ब्लैक और बाद में बर्फ़ ऊष्मामापी प्रयोग करने वाले लाव्वाज़िए तथा लाप्लास ने इसी प्रयोगात्मक सावधानी को अपनाया था।
\end{solution}
\end{exo}

\begin{exo}[खाद्य कैलोरी और उपापचयी शक्ति]{calorimetrie-calories-alimentaires}
\seoready{true}
\keywords{कैलोरी, किलोकैलोरी, खाद्य कैलोरी, यांत्रिक तुल्यांक, उपापचयी शक्ति, इकाइयाँ}

पोषण लेबल के अनुसार एक वयस्क औसतन प्रतिदिन लगभग $2000\,\text{kcal}$ लेता है। किलोकैलोरी पोषण में अभी प्रयुक्त गैर-SI इकाई है; $1\,\text{kcal}=4{,}18\times10^3\,\text{J}$।
\begin{enumerate}
  \item दैनिक ऊर्जा को जूल में बदलिए।
  \item ऊर्जा 24 घंटे में समान रूप से उपयोग होने पर औसत उपापचयी शक्ति वाट में निकालिए।
  \item ऊर्जा बार पर ``$210\,\text{kcal}$'' लिखा है। यदि यह ऊर्जा पूरी ऊष्मा बने तो प्रारम्भ में $20\,^\circ\text{C}$ वाले कितने पानी को $100\,^\circ\text{C}$ तक गर्म कर सकती है? $c_{\text{water}}=4{,}18\times10^3\,\text{J}\!\cdot\!\text{kg}^{-1}\!\cdot\!\text{K}^{-1}$ लीजिए।
  \item मानव शरीर वास्तव में इस ऊर्जा को किन रूपों में उपयोग करता है? गुणात्मक उत्तर दीजिए।
\end{enumerate}
\begin{indication}
प्रश्न 2 में $\mathcal P=E/\Delta t$। प्रश्न 3 में पहले जूल में बदलें, फिर $Q=mc\Delta T$ से $m$ निकालें।
\end{indication}
\begin{solution}
\textbf{1.}
\[
E=2000\times4{,}18\times10^3=8{,}36\times10^6\,\text{J}.
\]
\textbf{2.}
\[
\mathcal P=\frac{8{,}36\times10^6}{8{,}64\times10^4}\approx97\,\text{W}.
\]
पूरे दिन में मनुष्य की औसत शक्ति एक तापदीप्त बल्ब के लगभग बराबर है। यह प्रायः आश्चर्यजनक परिमाण-क्रम खाद्य कैलोरी को परिचित भौतिक इकाइयों में बदलने की उपयोगिता दिखाता है।

\textbf{3.}
\[
Q=210\times4{,}18\times10^3\approx8{,}78\times10^5\,\text{J}.
\]
$\Delta T=80\,\text{K}$ के लिए
\[
m=\frac{Q}{c\Delta T}
=\frac{8{,}78\times10^5}{4{,}18\times10^3\times80}
\approx2{,}63\,\text{kg}.
\]
अतः परिमाण-क्रम में एक ऊर्जा बार लगभग $2{,}5\,\text{kg}$ नल के पानी को उबालने लायक ऊर्जा रखता है।

\textbf{4.} शरीर ऊर्जा को गायब करने के अर्थ में ``खर्च'' नहीं करता, बल्कि पोषक तत्त्वों की रासायनिक ऊर्जा बदलता है। कुछ ऊर्जा पेशियों के काम, अंगों के संचालन, कोशिकीय विद्युत प्रवणताओं, अणु परिवहन और रासायनिक संश्लेषण में जाती है; कुछ ग्लाइकोजन या वसा के रूप में संचित हो सकती है। अन्ततः प्रयुक्त ऊर्जा का अधिकांश ऊष्मा बनकर क्षय होता है, शरीर का ताप बनाए रखने में सहायक होकर परिवेश में चला जाता है। थोड़ा भाग अपशिष्ट में बची रासायनिक ऊर्जा के रूप में शरीर से बाहर जाता है।
\end{solution}
\end{exo}
